% =============================================================================
\section{Background and Motivation}
\label{sec:bg}
% Budget: 2.0 pages. Draft FIRST (Step 5.1). Mixed audience: explain the OS
% mechanism for ML readers in under a page. Each observation = bold run-in head +
% measured evidence + "-> implication" used later by a design element.
% Close with research questions (CORE-style) that map to Section 5.
% =============================================================================

\subsection{Page Replacement, as the Kernel Sees It}
\label{sec:bg-mechanism}

\plan{For ML readers, in about 0.5 page. (1) A process has more virtual pages than its
resident limit allows; touching a non-resident page traps (a page fault); the kernel
must pick a resident \emph{victim}, write it back if it is dirty or has no swap copy,
and load the faulting page. (2) The decision: which resident page to evict; made once
per eviction, synchronously on the fault path; its budget is set by the fault it
serves (a disk read, often a write). (3) What the kernel can see: the MMU handles loads
and stores itself and only sets the accessed (A) and dirty (D) bits in the page-table
entry; the kernel reads and clears them when it scans. It never sees individual
accesses. Linux has the same limit for anonymous memory. (4) \clock and Aging as
heuristics over the A bit; \belady as the offline optimum (evict the page whose next
use is furthest away), optimal for faults but not for writes.}

\begin{figure}[t]
\centering
\figbox{1.55in}{The eviction path. Fault $\rightarrow$ resident limit reached
$\rightarrow$ \emph{choose victim} (policy: FIFO / \clock / Aging / LFU / \sys)
$\rightarrow$ write back if dirty or no swap copy $\rightarrow$ load page. Annotate what
the policy can read (A/D bits at scan time, per-page counters, swap-slot metadata) and
what it cannot (individual accesses).}{to draw (TikZ or vector); mechanism from
kernel/vmpage.c \texttt{reclaim\_frame}}
\caption{\plan{Takeaway caption: the policy decides once per eviction, on the fault path,
seeing only accessed/dirty bits.}}
\label{fig:mechanism}
\end{figure}

\subsection{Heuristics Leave Large Headroom}
\label{sec:bg-headroom}

\obs{1}{\clock is as good as exact LRU, and both are far from optimal.}
Across five workload families (\S\ref{sec:eval-setup}), \clock is within 1--2\% of
exact LRU on four, xv6's Aging is barely better than FIFO, and the only classical
policy that beats \clock anywhere, exact LFU, is catastrophic elsewhere
($\geq3\times$ \clock on graph, sort and matmul). \belady needs 43--54\% fewer faults
than \clock on the B-tree, key-value and graph workloads, 38\% on matmul, and only
10--17\% on sort (Figure~\ref{fig:headroom}).
$\rightarrow$ \emph{Recency and frequency both carry signal, but every fixed combination
fails on some workload; a policy should learn the combination.}
% src: ML_REPORT 6.1 classical_test / classical_heldout (simulator, geo-mean 5/10/20%)

\begin{figure}[t]
\centering
\figbox{1.6in}{Faults relative to \clock for FIFO, Aging, decayed LFU, LRU, exact LFU and
\belady, per workload family (test and held-out streams, geo-mean over 5/10/20\% memory).}
{exists as report/figures2/classical.pdf; re-plot without in-figure title, grayscale-safe}
\caption{\plan{Takeaway: \belady incurs 43--54\% fewer faults than \clock on btree, kv and
graph; no classical policy closes that gap.}}
\label{fig:headroom}
\end{figure}

\subsection{The Kernel's Coarse Signals Carry the Information}
\label{sec:bg-signals}

\obs{2}{What the kernel can sample nearly matches what an oracle trace records.}
Ranked within each eviction decision, the sampled access count \feat{sfreq} correlates
with the true next-use distance almost as strongly as the exact count \feat{freq}
($-0.45$ vs.\ $-0.52$ on btree, $-0.44$ vs.\ $-0.52$ on kv), and the refault distance
that Linux's workingset code already measures is as strong again ($+0.45$ btree,
$+0.36$ graph, $+0.35$ kv). Pooled correlation is misleading: recency's pooled
correlation on sort is $+0.87$, but within a decision it is $+0.05$
(Figure~\ref{fig:diagnostics}).
$\rightarrow$ \emph{Use kernel-observable features, and train and judge them on eviction
decisions, not on individual references} (\S\ref{sec:design-features},
\S\ref{sec:design-train}).
% src: ML_REPORT 6.2, report/results2/feature_diagnostics.csv

\begin{figure}[t]
\centering
\figbox{1.5in}{Within-decision rank correlation of each of the 12 features with the true
next-use distance, per workload (oracle tier vs.\ kernel tiers).}
{exists as report/figures2/diagnostics.pdf; re-plot for one column}
\caption{\plan{Takeaway: kernel-observable features are nearly as informative as oracle
features at the moment of the decision.}}
\label{fig:diagnostics}
\end{figure}

\subsection{A Naively Trained Policy Fails}
\label{sec:bg-naive}

\obs{3}{Trained on uniformly sampled candidates, a learned policy evicts the page the
program is streaming through.}
With \belady's choice plus 31 uniformly sampled candidates per recorded eviction, almost
every linear model was $\geq3\times$ \clock on PageRank, even one trained on PageRank
alone, and three rounds of DAgger~\citep{ross2011dagger} did not help. 97\% of its
victims had been loaded at the immediately preceding fault and were needed again a
median of 16 references later; \belady's choice at the same decisions had last been used
about 300 references earlier. Out of 167 resident pages, the one or two just touched are
almost never sampled, so the model never learned to keep them.
$\rightarrow$ \emph{The training distribution must contain just-used pages, and new pages
need protection} (\S\ref{sec:design-train}).
% src: ML_REPORT 7.1 (graph-pr2000x2-s2, validation)
\plan{This is the ``naive learned approach fails'' move reviewers reward (cf.\ Catur's
off-the-shelf RL losing to a rule). Consider a small figure: victim age distribution,
naive vs.\ \sys vs.\ \belady.} \todo{decide: figure or text only}

\subsection{Research Questions}
\label{sec:bg-rq}

Whether a learned policy is implementable in a kernel breaks into four questions,
answered in \S\ref{sec:eval}:
\begin{description}
\item[RQ1 Information.] Do kernel-observable signals suffice, or does a learned policy
  need information only a full trace has? (\S\ref{sec:eval-info})
\item[RQ2 Effect in the kernel.] Does the policy beat the kernel's best heuristic inside
  a real kernel, on unseen runs and unseen behaviours? (\S\ref{sec:eval-kernel})
\item[RQ3 Cost.] Can it be computed without floating point, and what does a decision
  cost on the eviction path? (\S\ref{sec:eval-cost})
\item[RQ4 Robustness.] What happens when the model is wrong, or meets a workload family
  it never saw? (\S\ref{sec:eval-robust})
\end{description}
